% Options for packages loaded elsewhere
\PassOptionsToPackage{unicode}{hyperref}
\PassOptionsToPackage{hyphens}{url}
\documentclass[
  12pt,
]{ctexart}
\usepackage{xcolor}
\usepackage{amsmath,amssymb}
\setcounter{secnumdepth}{-\maxdimen} % remove section numbering
\usepackage{iftex}
\ifPDFTeX
  \usepackage[T1]{fontenc}
  \usepackage[utf8]{inputenc}
  \usepackage{textcomp} % provide euro and other symbols
\else % if luatex or xetex
  \usepackage{unicode-math} % this also loads fontspec
  \defaultfontfeatures{Scale=MatchLowercase}
  \defaultfontfeatures[\rmfamily]{Ligatures=TeX,Scale=1}
\fi
\usepackage{lmodern}
\ifPDFTeX\else
  % xetex/luatex font selection
\fi
% Use upquote if available, for straight quotes in verbatim environments
\IfFileExists{upquote.sty}{\usepackage{upquote}}{}
\IfFileExists{microtype.sty}{% use microtype if available
  \usepackage[]{microtype}
  \UseMicrotypeSet[protrusion]{basicmath} % disable protrusion for tt fonts
}{}
\makeatletter
\@ifundefined{KOMAClassName}{% if non-KOMA class
  \IfFileExists{parskip.sty}{%
    \usepackage{parskip}
  }{% else
    \setlength{\parindent}{0pt}
    \setlength{\parskip}{6pt plus 2pt minus 1pt}}
}{% if KOMA class
  \KOMAoptions{parskip=half}}
\makeatother
\usepackage{longtable,booktabs,array}
\usepackage{calc} % for calculating minipage widths
% Correct order of tables after \paragraph or \subparagraph
\usepackage{etoolbox}
\makeatletter
\patchcmd\longtable{\par}{\if@noskipsec\mbox{}\fi\par}{}{}
\makeatother
% Allow footnotes in longtable head/foot
\IfFileExists{footnotehyper.sty}{\usepackage{footnotehyper}}{\usepackage{footnote}}
\makesavenoteenv{longtable}
\setlength{\emergencystretch}{3em} % prevent overfull lines
\providecommand{\tightlist}{%
  \setlength{\itemsep}{0pt}\setlength{\parskip}{0pt}}
\usepackage{bookmark}
\IfFileExists{xurl.sty}{\usepackage{xurl}}{} % add URL line breaks if available
\urlstyle{same}
\hypersetup{
  hidelinks,
  pdfcreator={LaTeX via pandoc}}

\author{SCX}
\date{2026}

\begin{document}

\section{SCX Project IP Attribution Statement}\label{scx-project-ip-attribution-statement}

\begin{quote}
This document records the intellectual property ownership and independence declaration of the SCX (State-Conditioned eXpertise) project. Last Updated: 2026-06-26
\end{quote}

\begin{center}\rule{0.5\linewidth}{0.5pt}\end{center}

\subsection{I. SCX Project Independence Declaration}\label{i-scx-project-independence-declaration}

\subsubsection{1.1 Core Nature: SCX is a Pure Math/ML Framework}\label{core-nature-scx-is-a-pure-mathml-framework}

SCX is a mathematical framework for \textbf{State-Conditioned eXpertise}. Its core contributions are:

\begin{itemize}
\tightlist
\item
  A definitional system (SCX Reliability, SCX Data Four-Classification, State-Wise Active Learning, Expert Routing, Noise vs Hard State distinction)
\item
  A set of formalized propositions (3 core theorems + 2 extended theorems)
\item
  Corresponding Python implementation
\end{itemize}

SCX \textbf{does not contain} any DFT (Density Functional Theory) computation code, VASP input files, or materials-science-specific computation logic.

\subsubsection{1.2 No Dependence on Institutional Resources}\label{no-dependence-on-institutional-resources}

\begin{longtable}[]{@{}
  >{\raggedright\arraybackslash}p{(\linewidth - 4\tabcolsep) * \real{0.2247}}
  >{\raggedright\arraybackslash}p{(\linewidth - 4\tabcolsep) * \real{0.1348}}
  >{\raggedright\arraybackslash}p{(\linewidth - 4\tabcolsep) * \real{0.6404}}@{}}
\toprule\noalign{}
\begin{minipage}[b]{\linewidth}\raggedright
Resource
\end{minipage} & \begin{minipage}[b]{\linewidth}\raggedright
Used?
\end{minipage} & \begin{minipage}[b]{\linewidth}\raggedright
Explanation
\end{minipage} \\
\midrule\noalign{}
\endhead
\bottomrule\noalign{}
\endlastfoot
Xiaogan Supercomputer (institutional equipment) & \textbf{No} & All SCX theoretical work, coding, and experiments completed on personal computer \\
Institutional network/storage & \textbf{No} & SCX project stored on personal computer and user's personal GitHub account \\
Institutional software licenses & \textbf{No} & SCX uses only open-source Python packages (numpy, scipy, scikit-learn, etc.) \\
Institutional personnel/students & \textbf{No} & SCX completed independently by a single person (the user) \\
\end{longtable}

\subsubsection{1.3 Experimental Independence}\label{experimental-independence}

SCX validation experiments include three categories:

\begin{enumerate}
\def\labelenumi{\arabic{enumi}.}
\tightlist
\item
  \textbf{Synthetic 2D data experiments}: Data self-generated by \texttt{experiments/synthetic/data\_generator.py}, no external data source needed
\item
  \textbf{MedMNIST / CIFAR}: Public academic benchmark datasets, freely downloadable by anyone
\item
  \textbf{MLIP case study}: Uses AlN v3 DFT data (EGP project output) as an optional, non-essential supplementary demonstration
\end{enumerate}

\begin{quote}
Synthetic data + public benchmark datasets are already sufficient to validate all core claims of SCX. The MLIP case study is purely an ``icing on the cake'' supplementary demonstration.
\end{quote}

\subsubsection{1.4 AI Code Generation Statement}\label{ai-code-generation-statement}

\textbf{All code} of the SCX Python package was generated by AI tools (AI programming tools) under the user's mathematical framework and design direction. The user is:

\begin{itemize}
\tightlist
\item
  \textbf{Original creator of the mathematical framework} (definitional system, propositions, proofs, algorithm design)
\item
  \textbf{Architecture designer} (module partitioning, interface design, overall architecture)
\item
  \textbf{Delegated code generator} (provided design specifications to AI; AI generated implementation)
\end{itemize}

\begin{center}\rule{0.5\linewidth}{0.5pt}\end{center}

\subsection{II. Boundary with EGP / AlGaN / DFT Projects}\label{ii-boundary-with-egp-algan-dft-projects}

\subsubsection{2.1 Two Independent Projects}\label{two-independent-projects}

\begin{longtable}[]{@{}lll@{}}
\toprule\noalign{}
Dimension & EGP Project (Paper 1-3) & SCX Project (Paper 4) \\
\midrule\noalign{}
\endhead
\bottomrule\noalign{}
\endlastfoot
Essence & MLIP (Machine Learning Interatomic Potential) & General ML theoretical framework \\
Dependency & Depends on DFT data (AlN v3) & Does not depend on any DFT data \\
Compute Resources & Uses Xiaogan Supercomputer & Personal computer suffices \\
Application Domain & Materials science & General machine learning \\
Code Repository & \texttt{egp/} & \texttt{SCX/} \\
\end{longtable}

\subsubsection{2.2 The Only Intersection: Conceptual Origin}\label{the-only-intersection-conceptual-origin}

The core concept of \textbf{``State-Conditioned eXpertise''} in SCX germinated from the gauge fixing work in EGP Paper 1 --- when splicing multiple ACE/PACE potentials, the user observed that the same potential function had different accuracy across different configuration regions.

However, \textbf{the SCX framework itself does not involve potential functions, DFT, or any materials science content at all}. It is a general ML theoretical framework that formalizes ``State-Conditioned eXpertise'' into a mathematical framework applicable to any supervised learning scenario (data valuation, active learning, expert routing, data compression).

\begin{quote}
\textbf{Analogy}: Newton conceived universal gravitation while observing a falling apple, but the law of universal gravitation is a universal physical theory independent of the specific variety of apple.
\end{quote}

\subsubsection{2.3 DFT Data Ownership of the EGP Project}\label{dft-data-ownership-of-the-egp-project}

AlN v3 DFT data was generated on the \textbf{Xiaogan Supercomputer} (institutional equipment) using VASP, and belongs to the EGP project output. SCX Project \textbf{does not claim independent ownership of this data}. The verifiability of the SCX framework does not depend on this data.

\begin{center}\rule{0.5\linewidth}{0.5pt}\end{center}

\subsection{III. Relationship with Institution/Research Group}\label{iii-relationship-with-institutionresearch-group}

\subsubsection{3.1 Severance Declaration}\label{severance-declaration}

\begin{itemize}
\tightlist
\item
  SCX Project \textbf{is not} a funded project of the user's institution or research group
\item
  SCX Project \textbf{did not use} the institution's research funding
\item
  SCX Project \textbf{did not use} the institution's computing resources, software licenses, or data
\item
  SCX Project \textbf{is not} part of the user's degree thesis or job-related research
\item
  SCX Project was completed \textbf{entirely on personal time, using personal equipment, personal computing resources}
\end{itemize}

\subsubsection{3.2 Potential Connections and Exclusions}\label{potential-connections-and-exclusions}

There is \textbf{no employment, funding, or supervisory relationship} between the SCX Project and the user's research group. Even if any general knowledge from the group existed (e.g., ML fundamentals), SCX's theoretical innovation --- the mathematical framework of State-Conditioned eXpertise --- is an original contribution independently proposed by the user on personal time.

\begin{center}\rule{0.5\linewidth}{0.5pt}\end{center}

\subsection{IV. Ownership Declaration Template}\label{iv-ownership-declaration-template}

\subsubsection{4.1 Ownership Attribution}\label{ownership-attribution}

All rights to the SCX Project (including but not limited to copyright, patent rights, attribution rights) belong to:

\begin{quote}
\textbf{[Reserved: Fill in rights holder information]}

Name: SCX
Date: \_\_\_\_\_\_\_\_\_\_\_\_\_2026.06.26\_\_\_\_\_\_\_\_\_\_\_\_\_
\end{quote}

\subsubsection{4.2 Code License}\label{code-license}

The release license for the SCX Python package is pending. Currently private repository, not distributed externally.

\subsubsection{4.3 Paper Authorship}\label{paper-authorship}

If academic papers are published, SCX authorship is determined by degree of independent contribution:

\begin{longtable}[]{@{}ll@{}}
\toprule\noalign{}
Contribution & Role \\
\midrule\noalign{}
\endhead
\bottomrule\noalign{}
\endlastfoot
Mathematical framework (definitions, propositions, proofs) & \textbf{User} (100\%) \\
Architecture design & \textbf{User} (100\%) \\
Code implementation & AI (AI programming tools) generated per user guidance \\
Experiment design and execution & \textbf{User} (100\%) \\
\end{longtable}

\begin{quote}
\textbf{AI tools are not entitled to authorship}. Per academic publishing norms (COPE, Nature, arXiv, etc.), AI tools shall not be listed as authors, but their use should be declared in the Acknowledgments or Methods section.
\end{quote}

\begin{center}\rule{0.5\linewidth}{0.5pt}\end{center}

\subsection{V. Legal and Ethical Declaration}\label{v-legal-and-ethical-declaration}

The content stated in this document truthfully reflects the development process, resource usage, and contribution attribution of the SCX Project. In case of dispute, the right to further proof is reserved through the following means:

\begin{enumerate}
\def\labelenumi{\arabic{enumi}.}
\tightlist
\item
  Git commit history (all commits from personal computer)
\item
  AI programming tools conversation records (evidence of AI-generated code)
\item
  Personal computer file metadata (creation/modification timestamps)
\item
  Development Log (\texttt{DEVELOPMENT\_LOG.md} in this repository)
\end{enumerate}

\begin{center}\rule{0.5\linewidth}{0.5pt}\end{center}

\emph{This document was authored by the user on 2026-06-26 to record the intellectual property status of the SCX Project.}

\end{document}
